\documentclass[11pt]{article}

\usepackage[margin=1in]{geometry}
\usepackage{amsmath,amssymb}
\usepackage{enumitem}

\newcommand{\LRF}{\mathrm{LRF}}
\newcommand{\OPT}{\mathrm{OPT}}
\newcommand{\maxi}{\max}
\newcommand{\mini}{\min}

\title{Maximums of $\lambda$ for $\alpha \ge \frac12$}
\author{}
\date{}

\begin{document}
\maketitle

We consider only instances with:
\begin{enumerate}[label=(\alph*)]
    \item $u < m$ long jobs, where $m$ is the number of machines, with processing time $p$.
    \item Additional short jobs have total workload $m$.
    \item $w_j = p_j$ for all $j$.
    \item $S_{\LRF}$ schedules all small jobs in $[0,1]$ and all long jobs start at $1$.
    \item $S_{\OPT}$ schedules all long jobs at time $0$, and all machines $i$ that execute a small job have the property
    \[
        C_{\max}^i(S_{\OPT}(I)) = C_{\max}^{\mini}(S_{\OPT}(I)).
    \]
\end{enumerate}

\section*{The Schedule $S_{\OPT}$}

For case (1) of Corollary 5,(REFERENCE MISSING)
\[
    C_{\max}^i(S_{\OPT}(I)) = C_{\max}^{\mini}(S_{\OPT}(I))
    \qquad \forall i.
\]
Moreover,
\[
    m + up \ge mp
    \quad\Longleftrightarrow\quad
    p \le \frac{1}{1-\frac{u}{m}}.
\]
Thus,
\begin{align*}
    C^*(I)(\alpha)
    &=
    \frac12 \left(
        \bigl(2\alpha-1\bigr)
        \sum_{j \in L(I)} p_j^2
        +
        \sum_{i=1}^m
        \bigl(C_{\max}^i(S_{\OPT}(I))\bigr)^2
    \right) \\
    &=
    \frac12 \left(
        \bigl(2\alpha-1\bigr)up^2
        +
        m\left(\frac{m+up}{m}\right)^2
    \right).
\end{align*}

\section*{The Schedule $S_{\LRF}$}

For $S_{\LRF}$,
\begin{align*}
    C_{\LRF}(I)(\alpha)
    &=
    \frac12 \left(
        \bigl(2\alpha-1\bigr)
        \sum_{j \in L(I)} p_j^2
        +
        \sum_{i=1}^m
        \bigl(C_{\max}^i(S_{\LRF}(I))\bigr)^2
    \right) \\
    &=
    \frac12 \left(
        \bigl(2\alpha-1\bigr)up^2
        +
        m-u
        +
        u(p+1)^2
    \right).
\end{align*}
% TODO: The handwritten blue note beside this expression appears to read:
% "Schweigel'schen does change of variables by Wissen of m."

Hence
\begin{align*}
    \lambda(u,p)
    &=
    \frac{
        \frac12\left(
            \bigl(2\alpha-1\bigr)up^2
            + m-u + u(p+1)^2
        \right)
    }{
        \frac12\left(
            \bigl(2\alpha-1\bigr)up^2
            + m\left(\frac{m+up}{m}\right)^2
        \right)
    } \\
    &=
    \frac{
        \bigl(2\alpha-1\bigr)up^2
        + m-u + up^2 + 2up + u
    }{
        \bigl(2\alpha-1\bigr)up^2
        + \frac{m^2+2mup+u^2p^2}{m}
    } \\
    &=
    \frac{
        2\alpha\frac{u}{m}p^2 + 1 + 2\frac{u}{m}p
    }{
        \bigl(2\alpha-1\bigr)\frac{u}{m}p^2
        + 1 + 2\frac{u}{m}p
        + \left(\frac{u}{m}\right)^2p^2
    }.
\end{align*}

This is for the case
\[
    \frac12 \le \alpha \le 1,
    \qquad
    0 \le \frac{u}{m} \le 1,
    \qquad
    0 \le p \le \frac{1}{1-\frac{u}{m}}.
\]
Wolfram Alpha gives the maximum under these constraints as
\[
    \frac54,
    \qquad \text{at } \alpha = \frac12,
\]
and
\[
    -
    \frac{
        2\left(
            4\sqrt{2}\sqrt{\frac{\alpha}{(2\alpha-1)^4}}\alpha^2
            - 2\alpha^2
            - 4\sqrt{2}\sqrt{\frac{\alpha}{(2\alpha-1)^4}}\alpha
            + \alpha
            + \sqrt{2}\sqrt{\frac{\alpha}{(2\alpha-1)^4}}
            - 1
        \right)
    }{
        (2\alpha-1)^2
    },
    \qquad
    \frac12 \le \alpha \le 1.
\]

Now for case (2) of Corollary 5:
\[
    C_{\max}^i(S_{\OPT}(I')) > C_{\max}^{\mini}(S_{\OPT}(I'))
    \qquad \text{for some } i.
\]
Again,
\[
    m + up \ge mp
    \quad\Longleftrightarrow\quad
    p \le \frac{1}{1-\frac{u}{m}}.
\]

For $S_{\OPT}$,
\begin{align*}
    C^*(I)(\alpha)
    &=
    \frac12 \left(
        \bigl(1 - 2(1-\alpha)\bigr)
        \sum_{j \in L(I)} p_j^2
        +
        \sum_{i=1}^m
        \bigl(C_{\max}^i(S(I))\bigr)^2
    \right) \\
    &=
    \frac12 \left(
        \bigl(1 - 2(1-\alpha)\bigr)up^2
        +
        (m-u)\left(\frac{m}{m-u}\right)^2
        +
        up^2
    \right) \\
    &=
    \frac12 \left(
        2\alpha up^2
        +
        (m-u)\left(\frac{m}{m-u}\right)^2
    \right).
\end{align*}

The value of $S_{\LRF}$ stays unchanged, so
\begin{align*}
    \lambda(I)
    &=
    \frac{
        \frac12\left(2\alpha up^2 + m + 2up\right)
    }{
        \frac12\left(2\alpha up^2 + (m-u)\left(\frac{m}{m-u}\right)^2\right)
    } \\
    &=
    \frac{
        2\alpha \frac{u}{m}p^2 + 1 + 2\frac{u}{m}p
    }{
        2\alpha \frac{u}{m}p^2 + \frac{m}{m-u}
    } \\
    &=
    \frac{
        2\alpha \frac{u}{m}p^2 + 1 + 2\frac{u}{m}p
    }{
        2\alpha \frac{u}{m}p^2 + \frac{1}{1-\frac{u}{m}}
    }.
\end{align*}

For
\[
    0 \le \frac{u}{m} \le 1,
    \qquad
    \frac12 \le \alpha \le 1,
    \qquad
    p > \frac{1}{1-\frac{u}{m}},
\]
Wolfram Alpha gives for this the maximum
\[
    \frac{1}{1-\frac{2}{(2\sqrt{\alpha}+\sqrt{2})^2}},
    \qquad \text{for } \alpha \ge \frac12,
\]
and the maximum is $-\infty$ for $\alpha \le \frac12$.

For
\[
    \frac{u}{m}
    =
    \frac{
        \sqrt{2}\sqrt{\frac{\alpha}{(\alpha-2)^2}}\alpha
        - 2\sqrt{2}\sqrt{\frac{\alpha}{(\alpha-2)^2}}
        + 2
    }{
        4 - 2\alpha
    },
\]
determine which case is bigger for which $\alpha$.




\end{document}
